% !TeX encoding = UTF-8 !TeX spellcheck = en_US
\documentclass[11pt, a4paper]{article}

% The title of the current document to be produced.
\newcommand{\doctitle}{Assignment \#4}
\newcommand{\assignments}{\bluetext{\textbf{Assignment:}}}
%
\setlength{\unitlength}{1in}
\renewcommand{\arraystretch}{2}
\setlength{\columnsep}{2cm}

\newcommand\dunderline[3][-1pt]{{%
  \sbox0{#3}%
  \ooalign{\copy0\cr\rule[\dimexpr#1-#2\relax]{\wd0}{#2}}}}

\usepackage[sfdefault]{atkinson}
\usepackage[T1]{fontenc}

\def\code#1{\bluetext{\texttt{#1}}}


%------------------------------------------------------------
% Import commands for both teacher and course information.  | NOTE: Change your
% teacher and course info in these files. |
%------>------>------>------>------>------>------>------>-->|
%
%------------------------------------------------------------
%-- Import packages and custom command definitons.          |
%------>------>------>------>------>------>------>------>-->|
%% ==================================
%% ===== Teacher info             ===
%% ==================================
%%
\newcommand{\instructor}{Katelyn Breivik}
\newcommand{\office}{Wean 8319}
\newcommand{\hours}{3pm on Mondays}
\newcommand{\phone}{}
\newcommand{\college}{}
\newcommand{\email}{kbreivik@andrew.cmu.edu}
\newcommand{\faculty}{}
\newcommand{\department}{Department of Physics}
                              %|
%% ==================================
%% ===== Course-specific commands ===
%% ==================================

%- Instructions: change course info here. 
\newcommand{\semester}{Fall 2023}

\newcommand{\ponderation}{2-4-3 (Theory-Lab-Homework)}
\newcommand{\coursetitle}{Astrophysics of Stars and the Galaxy}
\newcommand{\coursenumber}{[33-467]}
\newcommand{\prerequisite}{33-224, 33-234}
 
\input{../includes/packages-imports}                          %|  
\input{../includes/custom-commands}   
%
%---> Genereate & inject metadata                           |
\input{../includes/hyperef.doc.info}                          %|
%------------------------------------------------------------

\topmargin -70pt
\begin{document} 

\normalfont

%-------------------------------------------------------------
%-- Make the header of the document                          |
%------>------>------>------>------>------>------>------>--> |
%---------------------------------------------------------------------
%-- The following produces the document header including the title.  |
%---------------------------------------------------------------------

\noindent % <-- need to have this first.

\hfill	
\begin{minipage}{0.60\textwidth}%
	\raggedleft%
	{\Large \textsc{\doctitle}\par}
	\doublerule 
\end{minipage}%
\vspace{0.9cm}

  
%-------------------------------------------------------------
%-- Problem # 1                                              |
%------>------>------>------>------>------>------>------>--> |
\noindent \dunderline{1.0pt}{\textbf{Problem \#1: Radiation Transport (18 points)}}\\
The most important form of energy transport for stars on the main sequence
(which is $90\%$ of the stellar lifetime), is radiation transport.
\vspace{8pt}

\noindent \textbf{Problem \#1a: How old are those photons? (8 points)}\\
Even though we don't always get to see them through Pittsburgh's clouds, photons
are leaving the surface of the Sun all the time. Estimate the minimum time it
takes a photon to travel from the core of the Sun to the surface, assuming that
the photon retains it's identity every time it is absorbed and re-emmited. This
is probably not realistic, but I haven't yet had the opportunity to ask a photon
if it remembers all the places it has been absorbed and emitted.
\vspace{4pt}

\noindent ($i$) First, calculate the mean free path of a photon in the interior
regions of the Sun noting that the minimum opacity of fully ionized hydrogen is
set by the electron scattering cross section to be
$\kappa_{\rm{es}}=0.4\,\rm{cm^2/g}$. \emph{Hint:} You can use the average
density that you calculated in Homework 3 for a $1\,\rm{M_{\odot}}$ star. (2
points)
\vspace{4pt}

\noindent ($ii$) Next, using the speed of the photon, the mean free path, and
the radius of the sun, derive an equation for the time it takes for the photon
to travel from the center of the sun to its surface. I think the easiest way to
do this is using the dimensions of each variable, but you might think a
different way is easiest! (4 points)
\vspace{4pt}

\noindent ($iii$) How does your result compare to the thermal timescale of the
Sun? (2 points)
\vspace{8pt}

\noindent \textbf{Problem \#1b: Conceptual stuff (4 points)}\\
\vspace{2pt} 
Why does the diffusion approximation for radiation transport break down near the
surface of a star? Supply your answer in full at least one full sentence. 

\vspace{8pt}
\noindent \textbf{Problem \#1c: How small is a `small' temperature gradient? (6
points)}\\
Estimate the temperature gradient, $\partial T/\partial r$, in a gas blob inside
of the Sun where $T\sim10^7\,\rm{K}$, $r \sim 0.2\,\rm{R_{\odot}}$, $l(r) \sim
3.8\times10^{33}\,\rm{erg/s}$, $\kappa_{\rm{es}}\sim 0.4\,\rm{cm^2/g}$. What is
the change in the temperature over the range of a single mean free path of a
photon?

\vspace{10pt}
\noindent \dunderline{1.0pt}{\textbf{Problem \#2: Figures and captions (20
points)}}\\
\noindent Your project presentations are coming up soon! As part of your
presentation, you should plan to include at least one original figure that you
(and your partner if you have one) make. In the presentation, you'll be able to
walk the audience through the figure. This isn't the only way that figures are
presented though!\\

\noindent For this problem, you should create a figure (you and your partner can
submit the same figure) that you could use in your presentation. The figure
should illustrate some conclusion or result from your project and should have
adequate axis labels and a legend if necessary. Once you have a figure, you
should also write a figure caption that describes the necessary information for
a reader to understand and interpret the figure themselves. It's also a good
idea to add in a punchline that you want your reader to take away from the
figure. \emph{You must write your own figure caption even if you and your
partner share the figure.}\\

\noindent \emph{\large Hint!} If you need some inspiration, take a look at the
some figures in the papers you searched for on ads in homework 2.

\end{document} 
